% Prostudováno: ucebnice/02_Funkce_a_rovnice/Linearni_funkce_a_rovnice/Linerání rovnice.tex,
% část Lineární nerovnice o jedné neznámé (ř. 503–690); nejednotné zápisy intervalů ve zdroji byly ověřeny nezávisle.
% Realisticky.cz: externi_zdroje/realisticky/02_Funkce_a_rovnice/02_Linearni_rovnice_a_nerovnice/020205_Nerovnice_upravy_nerovnic/020205_Lekce.pdf
% a 020206_Linearni_nerovnice/020206_Lekce.pdf; využito vysvětlení obratu znaménka pomocí pořadí čísel a grafu. Úlohy vlastní.
\documentclass[11pt,a4paper]{article}
% Spolecny styl pro ucitelske pripravy MAT (Work/MAT)
% Verze 1.1, 2026-10-23
\NeedsTeXFormat{LaTeX2e}
\ProvidesPackage{priprava}[2026/10/23 v1.1 Sablona ucitelskych priprav MAT]

% Zakladni typografie a cestina
\RequirePackage[T1]{fontenc}
\RequirePackage[utf8]{inputenc}
\RequirePackage[czech]{babel}
\RequirePackage{lmodern}
\RequirePackage{microtype}
\RequirePackage{amsmath,amssymb}
\RequirePackage{xcolor}
\RequirePackage{enumitem}
\RequirePackage[most]{tcolorbox}
\RequirePackage{tabularx,booktabs,array}
\RequirePackage{fancyhdr}
\RequirePackage{needspace}
\RequirePackage{geometry}
\RequirePackage[hidelinks]{hyperref}

\geometry{a4paper,top=17mm,bottom=18mm,left=23mm,right=23mm}
\definecolor{notegray}{gray}{0.39}
\setlength{\parindent}{0pt}
\setlength{\parskip}{2.5pt plus 1pt minus 0.5pt}
\setlist[itemize]{topsep=2pt,itemsep=1pt,parsep=0pt,leftmargin=1.45em}
\setlist[enumerate]{topsep=2pt,itemsep=2pt,parsep=0pt,leftmargin=1.65em}
\setcounter{secnumdepth}{0}
\renewcommand{\arraystretch}{1.12}

% Zahlavi: vlevo rocnik | globalni tema, vpravo cislo: nazev hodiny
% Zapati zustava zcela prazdne.
\pagestyle{fancy}
\fancyhf{}
\renewcommand{\headrulewidth}{0.35pt}
\renewcommand{\footrulewidth}{0pt}
\setlength{\headheight}{14pt}
\newcommand{\nastavPripravu}[4]{%
  \fancyhead[L]{\small #1\enspace\textbar\enspace #2}%
  \fancyhead[R]{\small #3: #4}%
}

% Minimalisticke cernobile boxy. Kazda definice zvlast.
\tcbset{commonlesson/.style={enhanced,breakable=false,colback=white,colframe=black,
  colbacktitle=white,coltitle=black,boxrule=0.45pt,
  arc=0.8mm,outer arc=0.8mm,left=3mm,right=3mm,top=1.6mm,bottom=1.6mm,
  toptitle=1mm,bottomtitle=0.7mm,fonttitle=\bfseries,
  attach boxed title to top left={xshift=2.5mm,yshift=-1.15mm},
  boxed title style={colback=white,colframe=white,boxrule=0pt,left=1mm,right=1mm}}}
\newtcolorbox{zadaniBox}[1][]{commonlesson,title={Zadání},#1}
\newtcolorbox{prikladBox}[1][]{commonlesson,title={Příklad},#1}
\newtcolorbox{definiceBox}[1][]{commonlesson,title={Definice},#1}

% Spolecny box pro zadani a jeho reseni, bez mezery mezi castmi.
% Pouziti:
%   \begin{prikladReseniBox}
%     Zadani ulohy.
%     \tcblower
%     \begin{align*} ... \end{align*}
%   \end{prikladReseniBox}
% Spodni cast ma jemny sedy podklad a delici linku; lze v ni sazet
% vice radku matematiky. Obe casti maji jen jeden spolecny ramecek.
\tcbset{solutionlesson/.style={
  commonlesson,
  enhanced jigsaw,
  breakable,
  colbacklower=black!3,
  segmentation style={solid,draw=black!25,line width=0.35pt},
  fontlower=\small\color{notegray},
  before lower={\noindent\textbf{Řešení:}\enspace},
  leftlower=3mm,rightlower=3mm
}}
\newtcolorbox{prikladReseniBox}[1][]{solutionlesson,title={Příklad},#1}
\newtcolorbox{zadaniReseniBox}[1][]{solutionlesson,title={Zadání},#1}

% Pomucky pro autora pripravy: sedive poznamky; \reseni pro starsi pripravy
\newcommand{\poznamka}[1]{\par\begingroup\small\color{notegray}\itshape
  \textnormal{\textbf{Poznámka pro učitele.}}\enspace #1\par\endgroup}
\newcommand{\reseni}[1]{\par\begingroup\small\color{notegray}
  \textbf{Očekávané řešení:}\enspace #1\par\endgroup}
\newcommand{\blok}[1]{\needspace{5\baselineskip}\subsection*{#1}}
\newcommand{\vyrok}[1]{\uv{#1}}

\usepackage{tikz}
\nastavPripravu{1. ročník}{Rovnice a nerovnice}{75}{Lineární nerovnice}
\begin{document}
\blok{Proč se někdy mění směr nerovnosti}
Rovnice hledá rovnost, nerovnice hledá všechny argumenty, při nichž je jedna strana větší či menší. Například $2<5$, ale po násobení číslem $-1$ platí $-2>-5$. Záporné násobení převrátí pořadí čísel na ose. Přičtení stejného čísla pořadí nemění; násobení kladným číslem také ne.
\begin{definiceBox}
\textbf{Řešení nerovnice} $L(x)<P(x)$ na oboru $D$ je každé $x\in D$, pro které po dosazení platí nerovnost; obdobně pro $>,\le,\ge$. Množina všech řešení se obvykle zapisuje intervalem či sjednocením intervalů.
\end{definiceBox}
\begin{definiceBox}
\textbf{Lineární nerovnice o jedné neznámé} je nerovnice převoditelná na $ax+b\mathrel{\square}0$, kde $a,b\in\mathbb R$, $a\ne0$ a $\square$ je jeden ze znaků $<,>,\le,\ge$. Při $a=0$ rozhodneme pravdivost číselné nerovnosti zvlášť.
\end{definiceBox}
\begin{definiceBox}
\textbf{Ekvivalentní úpravy nerovnice}: k oběma stranám můžeme přičíst stejný výraz. Násobení nebo dělení kladným číslem zachová směr, záporným číslem jej obrátí. Nulou nedělíme; při násobení nulou by se ztratila informace o pořadí.
\end{definiceBox}
\begin{prikladReseniBox}
Řešte $-3x+4\le10$ a výsledek znázorněte na číselné ose. Ověřte jeden bod uvnitř a jeden mimo množinu řešení.
\tcblower
Odečteme $4$: $-3x\le6$. Dělíme $-3<0$, proto obrátíme směr: $x\ge-2$. Tedy $K=\langle-2;\infty)$. Kontrola: $x=0$ dává $4\le10$ (vyhovuje), $x=-3$ dává $13\le10$ (nevyhovuje). Bod $-2$ patří do řešení, protože původní nerovnost dává $10\le10$.
\end{prikladReseniBox}
\begin{center}
\begin{tikzpicture}[x=.8cm,>=stealth]
\draw[->] (-5,0)--(5,0) node[right] {$x$};
\foreach \x in {-4,-3,-2,-1,0,1,2,3,4}{\draw (\x,.09)--(\x,-.09) node[below] {\small $\x$};}
\draw[very thick] (-2,0)--(4.7,0);\fill (-2,0) circle (2.8pt);
\node[above] at (1.5,0) {\small $K=\langle-2;\infty)$};
\end{tikzpicture}
\end{center}
\begin{prikladReseniBox}
Řešte $2(x-1)>x+3$ a vysvětlete, proč mezní bod do výsledku nepatří.
\tcblower
$2x-2>x+3$, po odečtení $x$ a přičtení $2$ dostaneme $x>5$. Množina řešení je $(5;\infty)$. Pro $x=5$ mají obě strany hodnotu $8$, ale požadujeme ostrou nerovnost $8>8$, což neplatí. Např. $x=6$ dává $10>9$.
\end{prikladReseniBox}
\blok{Přepis na znaménko funkce}
Nerovnice $ax+b>0$ se ptá, na které části osy $x$ leží graf $y=ax+b$ nad osou $x$. Hranici nejprve najdeme z rovnice $ax+b=0$; potom podle znaménka směrnice či dosazení jednoho čísla rozhodneme stranu. Tento pohled vysvětluje interval, ale neomlouvá chybné algebraické obrácení nerovnosti.
\begin{prikladReseniBox}
Řešte $-2x+6>0$ početně i pomocí grafu přímky.
\tcblower
Početně $-2x>-6$, dělením $-2$ dostaneme $x<3$, tedy $K=(-\infty;3)$. Graf $y=-2x+6$ protíná osu $x$ v $(3;0)$ a vlevo od $3$ leží nad osou, protože klesá a například při $x=0$ má hodnotu $6>0$. Bod $3$ kvůli ostré nerovnosti nepatří.
\end{prikladReseniBox}
\blok{Samostatné procvičení}
\begin{zadaniReseniBox}
\begin{enumerate}[label=\alph*)]
\item Řešte $5-2x\ge1$.
\item Řešte $\frac{x-1}{3}<2$.
\item Řešte $3(x+1)\le x-5$ a znázorněte výsledek na ose.
\end{enumerate}
\tcblower
\begin{enumerate}[label=\alph*)]
\item $-2x\ge-4$, po dělení $-2$ je $x\le2$; $K=(-\infty;2\rangle$. Kontrola hranice: $5-4=1$.
\item Násobíme kladným $3$: $x-1<6$, tedy $x<7$ a $K=(-\infty;7)$; $x=7$ dává rovnost, ne ostré $<$.
\item $3x+3\le x-5$, tedy $2x\le-8$ a $x\le-4$; $K=(-\infty;-4\rangle$. Na ose plný bod $-4$ a polopřímka vlevo.
\end{enumerate}
\end{zadaniReseniBox}
\blok{Volitelné rozšíření a rezerva}
\begin{zadaniReseniBox}
Na oboru $x\in\mathbb N_0$ řešte $4x+7<30$ a určete největší přípustné řešení.
\tcblower
Na $\mathbb R$ vyjde $4x<23$, tedy $x<23/4=5{,}75$. Průnik s $\mathbb N_0$ je $K=\{0,1,2,3,4,5\}$, největší řešení je $5$. Kontrola: při $5$ je $27<30$, při $6$ je $31<30$ nepravda.
\end{zadaniReseniBox}
\begin{zadaniReseniBox}
Žák z $-4x<-8$ odvodil $x<2$. Určete chybu i správnou množinu řešení.
\tcblower
Při dělení záporným $-4$ musíme směr obrátit: $x>2$, $K=(2;\infty)$. Např. $x=3$ dává $-12<-8$; chybný návrh $x=1$ by dal $-4<-8$, což neplatí.
\end{zadaniReseniBox}
\blok{Domácí práce}
\begin{zadaniReseniBox}
\begin{enumerate}[label=\alph*)]
\item Řešte $7-3x>1$; rozhodněte o krajním bodu.
\item Řešte $\frac{2x+1}{4}\ge\frac{x-2}{2}$.
\item Vysvětlete na číslech, proč násobení záporným číslem obrací pořadí.
\end{enumerate}
\tcblower
\begin{enumerate}[label=\alph*)]
\item $-3x>-6$, tedy $x<2$ a $K=(-\infty;2)$. Při $x=2$ nastane rovnost $1=1$, bod nepatří.
\item Násobíme kladným $4$: $2x+1\ge2x-4$, tedy $1\ge-4$, což platí pro všechna $x$; $K=\mathbb R$.
\item Např. $1<3$, ale po násobení $-2$ platí $-2>-6$: obrazy čísel si na ose vymění pořadí.
\end{enumerate}
\end{zadaniReseniBox}
\poznamka{Jádro: význam a pravidlo 12 min, příklady 13 min, samostatné úlohy s kontrolou 17 min, organizace 3 min. Rezerva mimo součet; domácí práce asi 15 min. Při nejistotě nech ověřit mezní bod a jedno číslo na každé straně.}
\end{document}
